\documentclass[a4paper,oneside,12pt]{extarticle}

\usepackage[utf8]{inputenc}
\usepackage[english]{babel}
\usepackage{amssymb,amsthm,mathtools,enumitem,geometry,hyperref,booktabs}
\usepackage{graphicx}
\usepackage{xcolor}
\usepackage{tikz}
\usetikzlibrary{calc,patterns,decorations.pathreplacing}
\usepackage{tikz-cd}
\usepackage{pgfplots}
\usepgfplotslibrary{fillbetween}
\pgfplotsset{compat=1.18}
\usepackage{float}

\geometry{lmargin=15mm,rmargin=15mm,tmargin=10mm,bmargin=10mm,includeheadfoot}
\pagestyle{plain}
\frenchspacing

\theoremstyle{plain}
\newtheorem{proposition}{Proposition}
\newtheorem{lemma}[proposition]{Lemma}

\theoremstyle{definition}
\newtheorem{definition}[proposition]{Definition}
\newtheorem{example}[proposition]{Example}

\theoremstyle{remark}
\newtheorem{remark}[proposition]{Remark}

\newcommand*{\E}{\mathbb{E}}
\newcommand*{\Var}{\mathrm{Var}}
\newcommand*{\mP}{\mathbb{P}}
\newcommand*{\mR}{\mathbb{R}}
\newcommand*{\mbF}{\mathcal{F}}
\newcommand*{\ve}{\varepsilon}




\begin{document}
\begin{titlepage}
\begin{center}

{\footnotesize\scshape\color{gray}
Working Notes in Quantitative Risk and Modeling
}

\vfill

{\LARGE\bfseries Homoscedasticity and Its Testing:\\[4pt]
From Conditional Expectation to the Breusch--Pagan Test}

\vspace{36pt}

{\normalsize Boris Garbuzov}

\vfill

{\Large\bfseries Part~8~/~12}

\vspace{6pt}

{\Large\bfseries The Envelope Theorem\\and a Historical Note on Score vs LM}

\vfill

{\small July 2026}

\end{center}
\end{titlepage}

\setcounter{tocdepth}{2}
\tableofcontents
\newpage

\section{The Envelope Theorem and Its Use in the LM Test}
\label{sec:envelope_lm}

\subsection{The envelope theorem in the simplest possible case}
\label{subsec:envelope_simple}
%--------------------------------------------------------------------

Before using the envelope theorem to generalise the shadow-price
finding, we state and prove it in the most elementary setting:
an \emph{unconstrained} one-variable maximisation with a parameter.

\paragraph{Setup.}
Let $f(x,c)$ be a smooth function of two variables, $x$ the variable
being optimised and $c$ a parameter.  For each $c$, let
\[
  x^*(c) = \arg\max_x f(x,c), \qquad
  V(c) = f(x^*(c), c) = \max_x f(x,c)
\]
be the maximiser and the resulting maximal value (the
\emph{value function}).  Assume $x^*(c)$ is differentiable in $c$.

\paragraph{Question.} How does $V(c)$ change as $c$ changes?
Naively, $V$ depends on $c$ in two ways: directly through the
second argument of $f$, and indirectly because the maximiser
$x^*(c)$ itself shifts.  The envelope theorem says the second
channel contributes nothing.

\begin{quote}
\textbf{Envelope theorem (elementary form).}
\[
  V'(c) = \frac{\partial f}{\partial c}\Big(x^*(c), c\Big).
\]
The total derivative of $V$ equals the \emph{partial} derivative of
$f$ with respect to $c$, evaluated at the optimum --- with no
correction term for $dx^*/dc$.
\end{quote}

\begin{proof}
By the chain rule,
\[
  V'(c) = \frac{\partial f}{\partial x}\big(x^*(c),c\big)\cdot
  \frac{dx^*}{dc} + \frac{\partial f}{\partial c}\big(x^*(c),c\big).
\]
Since $x^*(c)$ maximises $f(\cdot,c)$, the first-order condition
gives $\partial f/\partial x\,(x^*(c),c) = 0$.  The first term
vanishes, leaving $V'(c) = \partial f/\partial c\,(x^*(c),c)$.
\end{proof}

\paragraph{Worked example.}
Let $f(x,c) = -(x-c)^2 + c$.  For fixed $c$, this is maximised at
$x^*(c) = c$, giving $V(c) = f(c,c) = c$.  Direct differentiation:
$V'(c) = 1$.

\begin{figure}[H]
  \centering
  \includegraphics[width=0.68\textwidth]{fig_env_parabolas2d.png}
  \caption{Three sections of $f(x,c)=-(x-c)^2+c$ at $c=0,1,2$
  (red, blue, green): $-(x-0)^2+0$, $-(x-1)^2+1$, $-(x-2)^2+2$.
  Each parabola peaks at $x^*=c$ with peak height $V(c)=c$:
  the peaks trace out the line $V(c)=c$ as $c$ varies, visibly
  rising by one unit for each unit step in $c$.}
\end{figure}

Now check via the envelope theorem without re-deriving $x^*(c)$:
$\partial f/\partial c = 2(x-c) + 1$.  Evaluated at $x = x^*(c) = c$:
$2(c-c)+1 = 1$.  Matches.  The term $2(x-c)$, which would carry the
effect of $x^*$ shifting, evaluates to exactly zero at the optimum
--- this is the first-order condition doing its job.

\paragraph{On the definition of $V$ and where it can fail.}
A precise statement of $V$ is worth pausing on:
\[
  V(c) = \max_x f(x,c) = f(x^*(c), c).
\]
If $x^*(c)$ is not defined at some $c=c_0$ (no maximiser exists, or
it is not unique and we have not specified a selection rule), then
$V(c_0)$ is simply not defined either, and the envelope-theorem
formula cannot be applied there.  The theorem is a statement about
differentiating $V$ \emph{where $V$ is well-defined and
differentiable} --- it presupposes, rather than establishes,
the existence and smoothness of $x^*(\cdot)$ near $c_0$.  In our
example this is not an issue: $f(x,c)=-(x-c)^2+c$ is a downward
parabola in $x$ for every $c$, so $x^*(c)=c$ exists and is unique
and smooth everywhere.

\paragraph{The geometry in 3-d.}
The same example lets us see the envelope theorem geometrically.

\begin{figure}[H]
  \centering
  \includegraphics[width=0.72\textwidth]{fig_env_surface3d.png}
  \caption{The surface $z=f(x,c)=-(x-c)^2+c$ over the $(x,c)$ plane.
  Every cross-section at fixed $c$ (a slice parallel to the $x$-axis)
  is a downward parabola; its peak lies along the ridge
  $x=c$.  The value function $V(c)$ is the height of this ridge,
  read off as the curve traced by the peaks as $c$ varies --- the
  ``envelope'' of the family of parabolic sections, which gives the
  theorem its name.}
\end{figure}

\begin{figure}[H]
  \centering
  \includegraphics[width=0.6\textwidth]{fig_env_levellines.png}
  \caption{Contour (level-line) plot of $f(x,c)=-(x-c)^2+c$ in the
  $(c,x)$ plane.  The ridge line $x=c$ (where each $c$-section
  attains its maximum) appears as the locus where the contours are
  tangent to vertical lines $c=\mathrm{const}$ --- i.e.\ where
  $\partial f/\partial x=0$.  Along this ridge, the contour values
  increase steadily with $c$, visibly confirming $V(c)=c$.}
\end{figure}

\paragraph{Interpretation.}
At the optimum, $f$ is locally flat in the $x$-direction (that is
what ``optimum'' means).  So a small perturbation of $c$ changes
$V$ exactly as if $x^*$ had stayed put: the indirect effect through
$x^*(c)$ is a second-order (negligible) correction, not a
first-order one.  This is the entire content of the theorem, and
it generalises unchanged to the constrained case via the
Lagrangian, as we use next.

%--------------------------------------------------------------------
\subsection{An operator viewpoint, and a graphical proof of the equality}
\label{subsec:envelope_operator}
%--------------------------------------------------------------------

\paragraph{Two natural operators.}
Define two operators acting on functions $f(x,c)$ of two variables:
\[
  D_c : f \mapsto \frac{\partial f}{\partial c}
  \qquad\text{(partial differentiation in $c$)},
  \qquad
  M_x : f \mapsto \max_x f(x,\cdot)
  \qquad\text{(maximise out $x$, leaving a function of $c$)}.
\]
Applied to $f(x,c)=-(x-c)^2+c$ in the two possible orders:
\[
  D_c\big(M_x f\big) = D_c\, V = V'(c) = 1,
  \qquad
  M_x\big(D_c f\big) = \max_x\big[\,2(x-c)+1\,\big] = +\infty
\]
(the second expression is unbounded in $x$, since $2(x-c)+1\to+\infty$
as $x\to\infty$).  So the two operators manifestly do \emph{not}
commute: $D_cM_x f\ne M_xD_c f$ in general --- the right-hand
expression need not even be finite.

\paragraph{What the envelope theorem actually says.}
The correct statement is not an order-of-operations identity
between $D_c$ and $M_x$ applied independently, but a much more
specific fact: $D_c$ applied \emph{after} $M_x$ equals $D_c f$
\emph{evaluated at the particular point $x=x^*(c)$ singled out by}
$M_x$ --- not maximised again, just evaluated there:
\[
  D_c\big(M_x f\big)(c)
  \;=\;
  \big(D_c f\big)\Big|_{x = x^*(c)},
\]
i.e.\ $V'(c) = \partial f/\partial c\,(x^*(c), c)$.  This is the
boxed formula from the previous subsection, restated in operator
language.  It holds not because $D_c$ and $M_x$ commute as abstract
operators, but because of the special structure at $x^*(c)$: the
chain-rule term that would obstruct commutativity,
$\partial f/\partial x \cdot dx^*/dc$, vanishes identically there
by the first-order condition.  ``Commutativity'' is the wrong frame;
``evaluation at the optimiser, not re-optimisation'' is the right
one.

\paragraph{The graphical equality: tangency, not coincidence.}
The cleanest way to see $V'(c) = \partial f/\partial c\,(x^*(c),c)$
geometrically is to compare $V(c)$ against a \emph{naive} slice of
$f$ at a fixed value of $x$ --- say $x=1$ --- treated as a function
of $c$ alone:
\[
  h(c) := f(1, c) = -(1-c)^2 + c.
\]
$h$ and $V$ are \emph{different} functions: $V(c)=c$ is the true
value function (re-optimising $x$ at every $c$), while $h(c)$ holds
$x$ fixed at $1$ and only lets $c$ move.  They agree only where the
fixed point $x=1$ happens to equal the true optimiser, i.e.\ at
$c=1$ (since $x^*(1)=1$).  At every other $c$, $x=1$ is not optimal,
so $h(c) < V(c)$ strictly.

\begin{figure}[H]
  \centering
  \includegraphics[width=0.62\textwidth]{fig_env_tangency_2d.png}
  \caption{The true value function $V(c)=c$ (black, straight line)
  against the naive fixed-$x$ slice $h(c)=f(1,c)=-(1-c)^2+c$ (green,
  dashed parabola).  The two curves touch only at $c=1$ (red dot),
  where $x=1$ coincides with the true optimiser $x^*(1)=1$.
  Crucially, they are not just equal in value there but
  \emph{tangent}: $V'(1) = h'(1) = 1$, since $h'(c) = 2(c-1)+1$ gives
  $h'(1)=1=V'(1)$.  Away from $c=1$, the slopes differ.  This is the
  geometric content of the envelope theorem: $\partial f/\partial c$
  at the optimiser reproduces not just the value but the
  \emph{slope} of $V$, at that one point only.}
\end{figure}

\begin{figure}[H]
  \centering
  \includegraphics[width=0.78\textwidth]{fig_env_tangency_3d.png}
  \caption{The same comparison lifted onto the surface
  $z=f(x,c)=-(x-c)^2+c$.  The black ridge is the true value function
  $V(c)$, traced by the vertex $x^*(c)=c$ of each cross-sectional
  parabola.  The green dashed curve is the fixed-$x=1$ slice
  $h(c)=f(1,c)$, lying \emph{on} the surface but off the ridge except
  at one point.  The red dot marks $(x,c)=(1,1)$, where the green
  curve touches the black ridge --- here $x=1$ happens to be the
  optimiser for $c=1$, so the two curves are tangent at this single
  point, with matching slope $1$, even though they diverge
  everywhere else.}
\end{figure}

\paragraph{Why tangency, not equality of functions, is the right
picture.}
This resolves a natural first guess: one might hope that $V(c)$ and
some fixed-$x$ slice $h(c)=f(x_0,c)$ are globally related (e.g.\
shifted versions of each other, as happens when two functions share
\emph{all} their partial derivatives).  They are not --- $V$ and $h$
are honestly different functions, equal and tangent at exactly one
point ($c$ such that $x^*(c)=x_0$) and diverging elsewhere.  The
envelope theorem is a statement about a single point of contact,
not a global relationship between two surfaces or two curves: it
identifies the slope of the ridge with the slope of \emph{the one
cross-sectional slice that happens to pass through the ridge at
that point}, and no further.

\paragraph{A clean operator formula.}
The discussion above shows $D_c$ and $M_x$ do not commute as
operators.  But there \emph{is} a clean operator identity hiding in
the envelope theorem, once we introduce two further operators
alongside $D_c$ and $M_x$, with care about their domains and the
types (signatures) of what goes in and what comes out.

\begin{itemize}
\item $D_c$: partial differentiation in the second slot.
  \[
    D_c : f(x,c) \;\longmapsto\; g(x,c) := \frac{\partial f}{\partial c}(x,c).
  \]
  Type: (function of two variables) $\to$ (function of two
  variables).  Natural domain: $f$ differentiable in its second
  argument.

\item $M_x$: maximise out the first slot.
  \[
    M_x : f(x,c) \;\longmapsto\; h(c) := \max_x f(x,c).
  \]
  Type: (function of two variables) $\to$ (function of one
  variable).  Natural domain: $f$ bounded above in $x$ for each $c$
  (so the maximum exists).

\item $A_x$: the maximiser (argmax) in the first slot.
  \[
    A_x : f(x,c) \;\longmapsto\; a(c) := \operatorname*{arg\,max}_x f(x,c).
  \]
  Type: (function of two variables) $\to$ (function of one
  variable).  Natural domain: a strict subset of $M_x$'s domain ---
  we need not only a finite maximum but a \emph{unique} maximiser,
  so that $a(c)$ is single-valued.

\item $S_x$: substitution into the first slot.
  \[
    S_x : \big(g(x,c),\, a(c)\big) \;\longmapsto\; k(c) := g\big(a(c),c\big).
  \]
  Type: (function of two variables, function of one variable) $\to$
  (function of one variable).  Domain: all pairs of the stated
  signatures (no further restriction); $S_x$ is just composition,
  plugging $a(c)$ into the first slot of $g$ and leaving $c$ to
  appear in both places.
\end{itemize}

With these four operators precisely typed, the envelope theorem
becomes the identity
\[
  \boxed{\,D_c\big(M_x f\big) \;=\; S_x\big(D_c f,\; A_x f\big).\,}
\]

\emph{Why this is exactly right.}  Unwind both sides as functions
of $c$:
\[
  \text{LHS:}\quad D_c(M_xf)(c) = \frac{d}{dc}\Big[\max_x f(x,c)\Big]
  = V'(c).
\]
\[
  \text{RHS:}\quad S_x(D_cf,\,A_xf)(c)
  = (D_cf)\big(A_xf(c),\,c\big)
  = \frac{\partial f}{\partial c}\big(x^*(c),\,c\big).
\]
These are literally the two sides of the boxed formula from
\S\ref{subsec:envelope_simple}: $V'(c) = \partial f/\partial c\,
(x^*(c),c)$.  So the formula is not a coincidental shuffling of
symbols --- it is a faithful, type-correct restatement of the
envelope theorem, with $A_x f$ (not $M_x f$) feeding into the first
slot of $D_c f$.  Note carefully that $f$ appears \emph{twice} on
the right: once differentiated ($D_cf$) and once maximised over via
its argmax ($A_xf$) --- these are two different transformations of
the same underlying $f$, combined only at the very last step by
$S_x$.  This is the precise sense in which Question~1's intuition
``it should be some kind of substitution into the derivative'' is
correct, once $A_x$ (not the cruder $M_x$) is used to supply the
point of substitution.

\begin{remark}
The earlier (incorrect) guess $M_x(D_cf)$ fails for two reasons
simultaneously: (i) it maximises $D_cf$ over $x$ at each $c$
\emph{independently} of where $f$ itself is maximised, so it
re-optimises rather than substitutes; and (ii) as the example
showed, this re-optimisation can diverge to $+\infty$ even when
$M_xf$ itself is perfectly well-defined.  The corrected formula
$S_x(D_cf, A_xf)$ avoids both problems: $A_xf$ locates the same
point $x^*(c)$ used to build $V$, and $S_x$ merely evaluates
$D_cf$ there --- no second optimisation is performed.
\end{remark}

%--------------------------------------------------------------------

We showed $\lambda^* = f'(\theta_0)$, i.e.\ the multiplier is the
derivative of the constrained optimum with respect to the
constraint level.  How far does this identification extend?

\paragraph{Setup: optimum as a function of the constraint level.}
Consider the family of problems indexed by $c$:
\[
  \text{(P$_c$):} \qquad \max_\theta f(\theta)
  \quad \text{subject to} \quad g(\theta) = c.
\]
For each $c$ in some neighbourhood, let $\theta^*(c)$ denote the
solution and define the \emph{value function}
\[
  V(c) = f(\theta^*(c)) = \max_{\theta:\,g(\theta)=c} f(\theta).
\]
The shadow price is $V'(c)$ --- the rate at which the optimal value
changes as the constraint level shifts.  In economics this is the
general object of interest: differentiate the optimised outcome
with respect to any parameter of the problem, not only with
respect to one written as a multiplier.

\paragraph{The Probabilist's derivation: $\theta^*(c)$ as a local inverse.}
Before invoking the envelope theorem directly, it is worth seeing
\emph{why} $\theta^*(c)$ is differentiable at all, following the Probabilist's
argument explicitly.  At the unconstrained optimum (no $c$ yet),
the first-order condition is
\[
  f'(\theta^*) - \lambda^* g'(\theta^*) = 0,
\]
which, together with the constraint $g(\theta^*)=c$, can be solved
implicitly for $(\theta^*,\lambda^*)$ as functions of $c$.  The Probabilist's
notes record this via the implicit function theorem applied
directly to the stationarity system: writing $\theta^*=\theta^*(c)$
and differentiating $g(\theta^*(c)) = c$ with respect to $c$,
\[
  g'(\theta^*)\cdot\frac{d\theta^*}{dc} = 1
  \implies
  \frac{d\theta^*}{dc} = \frac{1}{g'(\theta^*)}.
\]
Equivalently: locally near the current $c$, $\theta^* = g^{-1}(c)$
--- the constrained optimiser is (locally) the inverse function of
the constraint map itself, since $g$ is a bijection near a point
where $g'\ne 0$.  This is exactly the relation
$\lambda^* = f'(\theta^*)/g'(\theta^*)$ combined with
$\theta^*=g^{-1}(c)$ that appears in the handwritten notes.

Having established that $\theta^*(c)$ (and likewise $\lambda^*(c)$)
are well-defined differentiable functions of $c$ near the current
value, we can now form the composite
\[
  y(c) := f\big(\theta^*(c)\big),
\]
and ask for $y'(c)$.  By the chain rule,
\[
  y'(c) = f'\big(\theta^*(c)\big)\cdot \frac{d\theta^*}{dc}(c).
\]
This is precisely $V(c)$ from the paragraph above ($y\equiv V$),
and substituting $d\theta^*/dc = 1/g'(\theta^*)$ recovers
\[
  V'(c) = \frac{f'(\theta^*(c))}{g'(\theta^*(c))} = \lambda^*(c),
\]
matching the boxed formula for $\lambda^*$ found earlier by direct
substitution into the stationarity conditions.  So the Probabilist's route
(differentiate the implicit system directly) and the envelope-theorem
route (below) arrive at the same conclusion; the envelope theorem
simply packages the cancellation that happens inside this
chain-rule computation, without requiring us to solve for
$d\theta^*/dc$ explicitly.

\paragraph{Recovering $\lambda$ via the envelope theorem.}
The elementary envelope theorem above had a single unconstrained
variable $x$.  To apply it to the Lagrangian, we enlarge the
variable solved for via the first-order condition from $\theta$
alone to the pair $(\theta,\lambda)$, and treat the joint
stationarity of $\mathcal{L}$ in this pair as the elementary
theorem's stationary point $x^*$ --- \emph{not} as a maximisation:
as the determinant computation in Part~5b-a, \S2.9 shows, $\mathcal{L}$
has a genuine saddle in $(\theta,\lambda)$, and for fixed $\theta\ne
c$ it is in fact \emph{unbounded} in $\lambda$ (linear, with no
extremum at all).  So there is no optimisation over $\lambda$ in any
ordinary sense --- the $\lambda$ first-order condition
$\partial\mathcal{L}/\partial\lambda=0$ exists only to \emph{recover
the constraint}, not to optimise anything.  This is legitimate because at a stationary point of a Lagrangian,
\emph{both} partial derivatives vanish --- exactly the single
first-order condition $\partial f/\partial x=0$ used in the proof
above, just written out in two components instead of one.

Form the Lagrangian for problem (P$_c$):
\[
  \mathcal{L}(\theta,\lambda;c) = f(\theta) - \lambda\,(g(\theta)-c).
\]
At the solution, $\theta^*(c)$ and $\lambda^*(c)$ are both functions
of $c$ (assumed differentiable, at least locally near the current
$c$).  Differentiating $V(c) = \mathcal{L}(\theta^*(c),\lambda^*(c);c)$
with respect to $c$ and using the envelope theorem (the
first-order conditions $\partial_\theta\mathcal{L}=0$ and
$\partial_\lambda\mathcal{L}=0$ kill the terms coming from
$d\theta^*/dc$ and $d\lambda^*/dc$, exactly as the single term
$\partial f/\partial x\cdot dx^*/dc$ vanished in the elementary
case):
\[
  V'(c) = \frac{\partial\mathcal{L}}{\partial c}
  \bigg|_{\theta^*(c),\,\lambda^*(c)}
  = \lambda^*(c).
\]
So $\lambda^*$ \emph{is} $V'(c)$: the Lagrange multiplier is the
derivative of the optimum with respect to the right-hand side of
the constraint.  This recovers our earlier finding
$\lambda^*=f'(\theta_0)$ as the special case $g(\theta)=\theta$,

$c=\theta_0$, where $V(c) = f(c)$ trivially and $V'(c)=f'(c)$.

\paragraph{How general is this?}
The derivation above did not use any special form of $g$ --- it
holds for any smooth equality constraint $g(\theta)=c$, scalar or
vector $\theta$, as long as the envelope theorem's regularity
conditions hold (differentiability of $\theta^*(\cdot)$,
$\lambda^*(\cdot)$ near $c$, and the constraint qualification).
So within the class of problems written as ``maximise $f$ subject
to $g(\theta)=c$,'' the shadow price of relaxing $c$ is always
exactly $\lambda^*$ --- one of the Lagrange multipliers (the one
attached to that constraint, if there are several).

\paragraph{Beyond constraint levels: shadow prices in general.}
In economics, ``shadow price'' is used more broadly: one
differentiates the optimised value with respect to \emph{any}
parameter of the problem, not only a constraint level appearing
on the right-hand side.  For example, if $f$ itself depends on a
parameter $a$ (say $f(\theta;a)$, as in a revenue function with a
price parameter $a$), then $\partial V/\partial a$ is also called
a shadow price, but it need not equal any Lagrange multiplier ---
the envelope theorem gives $\partial V/\partial a
= \partial f/\partial a$ evaluated at $\theta^*$, with no $\lambda$
involved, because $a$ does not appear in the constraint at all.

\paragraph{A further generalisation: parametrising $g$ itself.}
The Probabilist's suggestion: under suitable regularity, one can go further
and let the constraint be written as $g(\theta,c)=0$ directly ---
i.e.\ $c$ enters the left-hand side, possibly nonlinearly, rather
than appearing as a separate right-hand-side level subtracted off
at the end.  Write the family of problems as
\[
  \text{(P$_c$):} \qquad \max_\theta f(\theta)
  \quad\text{subject to}\quad g(\theta,c) = 0.
\]
This is strictly more general than (P$_c$) above: it reduces to it
when $g(\theta,c) := g(\theta) - c$, but allows $c$ to enter
nonlinearly, or even to enter both $\theta$ and $c$ jointly in ways
that do not separate.

Assume, for each $c$, a unique maximiser $\theta^*(c)$ exists and is
differentiable in $c$ (if several local maxima exist, the same
construction applies separately to each differentiable branch
$\theta^*_i(c)$; we work with one branch and assume uniqueness for
simplicity).  Define the value function
$f^*(c) := f(\theta^*(c))$, which is differentiable as a composition
of differentiable functions.  We re-derive the shadow-price formula
by the Probabilist's direct route (differentiating the constraint), and then
recover it again via the envelope theorem, to confirm the two routes
agree.

\begin{remark}[Branches and the manifold picture]
The case of several local maxima, handled above by picking a
differentiable branch $\theta^*_i(c)$, is the same phenomenon as
Question~2's restricted function $\tilde f$ on the constraint
manifold $\{g(\theta,c)=0\}$ for fixed $c$: each branch
$\theta^*_i(c)$ traces a local maximum of $\tilde f$ on that
manifold, and as $c$ varies these local maxima move continuously
along it.  Making this correspondence precise --- describing the
branches $\theta^*_i(c)$ intrinsically as critical points of
$\tilde f$ varying with $c$, rather than extrinsically via the
first-order condition in the ambient space --- requires the
implicit function theorem applied jointly in $(\theta,c)$, which is
a genuinely more delicate tool than anything used here. We do not
pursue this; for our purposes the extrinsic description (the
first-order condition $\nabla_\theta\mathcal{L}=0$ in the full
$\theta$-space, for each fixed $c$) is sufficient and is what the
derivation below uses throughout.
\end{remark}

\emph{Step 1: differentiate the constraint.}
Since $g(\theta^*(c), c) \equiv 0$ for every $c$, differentiate
both sides with respect to $c$ using the chain rule:
\[
  \frac{\partial g}{\partial\theta}\Big(\theta^*(c),c\Big)\cdot
  \frac{d\theta^*}{dc}
  +
  \frac{\partial g}{\partial c}\Big(\theta^*(c),c\Big) = 0.
\]
Solve for $d\theta^*/dc$ (valid whenever $\partial g/\partial\theta
\ne 0$, the constraint qualification):
\[
  \frac{d\theta^*}{dc} =
  -\,\frac{\partial g/\partial c}{\partial g/\partial\theta}
  \bigg|_{(\theta^*(c),\,c)}.
\]

\emph{Step 2: differentiate the value function.}
By the chain rule, $f^{*\prime}(c) = f'(\theta^*(c))\cdot
d\theta^*/dc$.  Substituting the result of Step~1:
\[
  f^{*\prime}(c) =
  -\,f'(\theta^*(c))\cdot
  \frac{\partial g/\partial c}{\partial g/\partial\theta}
  \bigg|_{(\theta^*(c),c)}.
\]

\emph{Step 3: bring in $\lambda^*$.}
From the Lagrangian first-order condition
$f'(\theta^*) - \lambda^*\,\partial g/\partial\theta\,(\theta^*,c) = 0$,
we have $f'(\theta^*) = \lambda^*\cdot \partial g/\partial\theta\,(\theta^*,c)$.
Substituting into the Step~2 formula, the
$\partial g/\partial\theta$ factors cancel:
\[
  \boxed{\,f^{*\prime}(c) = V'(c) =
  -\,\lambda^*(c)\,\frac{\partial g}{\partial c}\Big(\theta^*(c),c\Big).\,}
\]
Before this cancellation we had a ratio of two partial derivatives
of $g$; after substituting the Lagrangian condition we are left
with just the numerator term, multiplied by $\lambda^*$.

\emph{Recovering the simpler case.}
To confirm this generalises the earlier finding, set
$g(\theta,c) = \theta - c$ (our original convention).  Then
$\partial g/\partial c = -1$, and the boxed formula gives
$V'(c) = -\lambda^*\cdot(-1) = \lambda^*$ --- exactly the earlier
result $\lambda^*=f'(\theta_0)$.  So the simple case is the special
instance where $c$ enters $g$ additively with coefficient $-1$.

\paragraph{Worked example: the multiplier is not parametrisation-invariant.}
Take $f(\theta) = -(\theta-3)^2$ (unconstrained maximum at
$\theta=3$) and compare two constraints with the \emph{same} zero
set $\{\theta = c\}$ (for $\theta,c>0$), written two different ways:
\[
  g_1(\theta,c) = \theta - c, \qquad
  g_2(\theta,c) = \theta^2 - c^2.
\]
Both vanish exactly when $\theta=c$, so both give the same
$\theta^*(c) = c$ and the same value function
$V(c) = f(c) = -(c-3)^2$, hence the same true shadow price
$V'(c) = 6-2c$ (computed directly, with no reference to either
Lagrangian).

But the two multipliers differ:
\[
  \lambda_1^*(c) = \frac{f'(\theta^*)}{\partial g_1/\partial\theta}
  = \frac{6-2c}{1} = 6-2c,
  \qquad
  \lambda_2^*(c) = \frac{f'(\theta^*)}{\partial g_2/\partial\theta}
  = \frac{6-2c}{2c} = \frac{3-c}{c}.
\]
Checking the boxed formula for each: for $g_1$,
$\partial g_1/\partial c = -1$, so
$V'(c) = -\lambda_1^*\cdot(-1) = 6-2c$ --- correct.
For $g_2$, $\partial g_2/\partial c = -2c$, so
$V'(c) = -\lambda_2^*\cdot(-2c) = \frac{3-c}{c}\cdot 2c = 6-2c$
--- also correct, by a different route.  Both recover the same true
shadow price, but $\lambda_1^*\ne\lambda_2^*$ for $c\ne 0,3$: the
Lagrange multiplier itself depends on how the constraint is
\emph{written}, even though the constraint set and the shadow price
do not.  Only $V'(c)$ is an invariant of the problem; $\lambda^*$ is
an artefact of the chosen parametrisation of $g$.

\paragraph{Connection to the Wikipedia formulation and to inequality
constraints.}
The Wikipedia article on the envelope theorem treats this in
slightly more general terms still: it allows $f$ itself to depend
on $c$ (write $f(\theta,c)$, not just $f(\theta)$), and states the
result for inequality-constrained problems, which requires
Karush--Kuhn--Tucker machinery beyond what we need here.  Stripped
down to our equality-constraint, $f(\theta)$-independent-of-$c$
case, and writing $V(c)$ in place of Wikipedia's value-function
notation (matching our $f^*(c)$), their general statement
specialises to exactly the boxed formula above.  The cleanest way
to state it, avoiding the need to re-derive Steps~1--3 each time:
extend the Lagrangian itself to carry the $c$-dependence,
\[
  \mathcal{L}(\theta,\lambda;c) = f(\theta) - \lambda\, g(\theta,c),
\]
and apply the elementary envelope theorem of
\S\ref{subsec:envelope_simple} directly to this $\mathcal{L}$,
treating $(\theta,\lambda)$ jointly as the stationary point ``$x^*$''
of the elementary theorem (not as a maximiser):
\[
  V'(c) = \frac{\partial\mathcal{L}}{\partial c}
  \bigg|_{\theta^*(c),\lambda^*(c)}
  = -\lambda^*(c)\,\frac{\partial g}{\partial c}(\theta^*(c),c),
\]
identical to the boxed result, obtained with no need to separately
differentiate the constraint.  This is the version that should be
used in practice: state the parametrised Lagrangian once, then
differentiate it in $c$ at the optimum --- the envelope theorem
guarantees the $d\theta^*/dc$ and $d\lambda^*/dc$ terms vanish
automatically, exactly as they did in the elementary case.

%--------------------------------------------------------------------
\subsection{Building the chain rule graphically: the Probabilist's step-by-step
picture series}
\label{subsec:envelope_chainrule_pictures}
%--------------------------------------------------------------------

The envelope-theorem discussion above used the chain rule in
passing (\S\ref{subsec:envelope_simple}), but only for a single
variable $x$.  In the Lagrangian setting, and in general, we need
the chain rule for a function of \emph{two} variables $f(x,y)$
composed with a parametric path $x=x(c)$, $y=y(c)$.  Before
specialising to the envelope theorem (where one of the two
variables will turn out to be the identity, $y(c)=c$), the Probabilist's
notes build up the general picture in small steps, each one
illustrated separately.  We reproduce the series here exactly,
with commentary, and then continue it to the dot-product formula
it is building toward.

\paragraph{Setup.}
We have a surface $z=f(x,y)$ (Figure~\ref{fig:vova-env-1}) and a
path in the $(x,y)$-plane given parametrically by $c\mapsto
(x(c),y(c))$ (Figure~\ref{fig:vova-env-2}).  Importantly, $y$ is
{not} assumed to equal $c$ here --- both $x(c)$ and $y(c)$ are, for
now, general functions of the parameter $c$.  This generality
matters: in the envelope-theorem application later, $y(c)=c$ will
be a special (degenerate) case, and seeing the general chain rule
first prevents confusing the special structure of that case with
the general mechanism.

\begin{figure}[H]
  \centering
  \includegraphics[width=0.5\textwidth]{v1.png}
  \caption{The surface $z=f(x,y)$.  A generic dome shape, no
  parametrisation yet.}
  \label{fig:vova-env-1}
\end{figure}

\begin{figure}[H]
  \centering
  \includegraphics[width=0.78\textwidth]{v2.png}
  \caption{Left: the parameter axis $c$ (green), with the direction
  of increasing $c$ marked.  Right: the image of the map
  $c\mapsto(x(c),y(c))$ in the $(x,y)$-plane (red curve, arrows
  showing the direction of increasing $c$ along the path).}
  \label{fig:vova-env-2}
\end{figure}

\paragraph{The path as a curve in $(c,x,y)$-space, and its three
shadows.}
To keep track of both component functions $x(c)$ and $y(c)$
simultaneously, the Probabilist lifts the picture into the three-dimensional
space with axes $(c,x,y)$ --- note this is \emph{not} the
$(x,y,z)$ space of the surface $f$; $z=f$ has not entered yet.  The
red space curve is $c\mapsto(c,\,x(c),\,y(c))$, and it casts three
shadows (orthogonal projections) onto the three coordinate planes.

\begin{figure}[H]
  \centering
  \includegraphics[width=0.62\textwidth]{v3.png}
  \caption{The space curve $(c,x(c),y(c))$ in red, together with its
  shadow (orange) on the $(x,y)$-plane --- this orange shadow is
  exactly the red curve of Figure~\ref{fig:vova-env-2}, now seen as
  a projection rather than drawn directly.}
  \label{fig:vova-env-3}
\end{figure}

\begin{figure}[H]
  \centering
  \includegraphics[width=0.62\textwidth]{v4.png}
  \caption{A second shadow added: the magenta projection onto the
  $(c,x)$-plane is the graph of $x=x(c)$ --- the first component
  function on its own, as an ordinary function of $c$.}
  \label{fig:vova-env-4}
\end{figure}

\begin{figure}[H]
  \centering
  \includegraphics[width=0.62\textwidth]{v5.png}
  \caption{The third shadow added: the green projection onto the
  $(c,y)$-plane is the graph of $y=y(c)$.  All three shadows are now
  visible together: $(x,y)$ (orange, the path itself), $(c,x)$
  (magenta, the graph of $x(\cdot)$), $(c,y)$ (green, the graph of
  $y(\cdot)$).  This completes the description of the parametric
  path and its two component functions.}
  \label{fig:vova-env-5}
\end{figure}

\paragraph{The same construction for the derivatives.}
The Probabilist repeats the identical construction one level up: instead of
the curve $(c,x(c),y(c))$, take the curve of \emph{derivatives}
$(c,\,x'(c),\,y'(c))$, with primed axes $x',y'$.  The logic and the
three shadows are exactly analogous --- only the object being
plotted has changed, from positions to velocities.

\begin{figure}[H]
  \centering
  \includegraphics[width=0.62\textwidth]{v6.png}
  \caption{The same three-shadow construction, now for the
  derivative curve $(c,x'(c),y'(c))$: green is the $(c,y')$-shadow
  (graph of $y'(\cdot)$), magenta is the $(c,x')$-shadow (graph of
  $x'(\cdot)$).}
  \label{fig:vova-env-6}
\end{figure}

\begin{figure}[H]
  \centering
  \includegraphics[width=0.62\textwidth]{v7.png}
  \caption{The fourth projection added, in light blue: the shadow
  onto the $(x',y')$-plane itself.  This is the one we did not draw
  in the first (position) series --- it is the new, essential
  object: the image of the velocity vector $(x'(c),y'(c))$, traced
  out as $c$ varies.}
  \label{fig:vova-env-7}
\end{figure}

\paragraph{Isolating the velocity vector.}
The $(x',y')$-shadow is now extracted on its own and reinterpreted:
rather than viewing it as a curve traced by a moving point, each
point $(x'(c_0),y'(c_0))$ on it is drawn as a vector emanating from
the origin --- the instantaneous \emph{velocity} of the path
$(x(c),y(c))$ at parameter value $c_0$.

\begin{figure}[H]
  \centering
  \includegraphics[width=0.5\textwidth]{v8.png}
  \caption{The point $(x'(c),y'(c))$ on the $(x',y')$-plane, shown as
  a radius vector from the origin.  This is the velocity vector of
  the path at the corresponding parameter value --- the object that
  will be dotted with the gradient of $f$ below.}
  \label{fig:vova-env-8}
\end{figure}

\paragraph{Returning to the surface: lifting the path onto $f$.}
The second half of the series returns to the surface $z=f(x,y)$
and the floor path $(x(c),y(c))$, now shown together and then
combined: the floor path is lifted vertically onto the surface,
producing the composite curve
$c\mapsto\big(x(c),\,y(c),\,f(x(c),y(c))\big)$.

\begin{figure}[H]
  \centering
  \includegraphics[width=0.45\textwidth]{v9.png}
  \caption{The surface $z=f(x,y)$, redrawn and explicitly labelled.}
  \label{fig:vova-env-9}
\end{figure}

\begin{figure}[H]
  \centering
  \includegraphics[width=0.55\textwidth]{v10.png}
  \caption{The floor path $\{x=x(c),\,y=y(c)\}$, redrawn in
  isolation, matching Figure~\ref{fig:vova-env-2}.}
  \label{fig:vova-env-10}
\end{figure}

\begin{figure}[H]
  \centering
  \includegraphics[width=0.5\textwidth]{v11.png}
  \caption{Surface and floor path shown together, the path still
  lying flat on the $z=0$ plane beneath the surface --- not yet
  lifted onto it.}
  \label{fig:vova-env-11}
\end{figure}

\begin{figure}[H]
  \centering
  \includegraphics[width=0.62\textwidth]{v12.png}
  \caption{The floor path lifted vertically onto the surface,
  producing the red space curve
  $\{x=x(c),\,y=y(c),\,z=f(x,y)\}$ --- i.e.\ the graph of
  $h(c):=f(x(c),y(c))$ traced out on the surface itself.  This is
  the final object of the Probabilist's series: the composite function whose
  derivative we want, displayed not as an abstract formula but as
  an actual curve riding on the surface, directly above its shadow
  on the floor.}
  \label{fig:vova-env-12}
\end{figure}

\paragraph{Continuing the series: assembling the chain rule.}
With every piece now in view --- the surface, the floor path, the
lifted curve $h(c)$, and the velocity vector $(x'(c),y'(c))$ --- we
can assemble them into the chain-rule formula.  We continue the Probabilist's
series with three more pictures, completing the construction he
sketched and naming as the goal: \emph{this is a dot product of two
vectors.}

\begin{figure}[H]
  \centering
  \includegraphics[width=0.68\textwidth]{v13_lifted_curve.png}
  \caption{All the pieces together for a concrete example
  $f(x,y)=4-x^2-0.6y^2$ with path $x(c)=c$, $y(c)=\tfrac12 c^2$
  (chosen so that $y$ is \emph{not} the identity, matching the Probabilist's
  emphasis that $y$ need not equal $c$): the orange floor curve
  $(x(c),y(c))$, and the red lifted curve $h(c)=f(x(c),y(c))$ riding
  on the surface directly above it.  Dotted vertical segments show
  the lift at a few sample parameter values.  This is the direct
  continuation of Figure~\ref{fig:vova-env-12} to a fully worked
  numerical case.}
  \label{fig:vova-env-13}
\end{figure}

\paragraph{The chain rule itself.}
$h(c) = f(x(c), y(c))$ is a composition of $f:\mathbb{R}^2\to
\mathbb{R}$ with the path $c\mapsto(x(c),y(c)):\mathbb{R}\to
\mathbb{R}^2$.  The multivariate chain rule gives
\[
  h'(c) = \frac{\partial f}{\partial x}\big(x(c),y(c)\big)\cdot x'(c)
  \;+\;
  \frac{\partial f}{\partial y}\big(x(c),y(c)\big)\cdot y'(c).
\]
Writing the two partial derivatives as the gradient vector
$\nabla f = (\partial f/\partial x,\ \partial f/\partial y)$,
evaluated at $(x(c),y(c))$, and the two derivatives $x'(c),y'(c)$
as the velocity vector from Figure~\ref{fig:vova-env-8}, the
right-hand side is exactly a dot product:
\[
  \boxed{\,h'(c) = \nabla f\big(x(c),y(c)\big)\,\cdot\,
  \big(x'(c),\,y'(c)\big).\,}
\]

\begin{figure}[H]
  \centering
  \includegraphics[width=0.55\textwidth]{v14_velocity_components.png}
  \caption{The velocity vector $(x'(c),y'(c))$ (blue) decomposed
  into its two axis components $x'(c)$ (orange, along the
  $x'$-axis) and $y'(c)$ (green, along the $y'$-axis) --- the two
  terms that appear, each multiplying a partial derivative of $f$,
  in the chain-rule sum above.}
  \label{fig:vova-env-14}
\end{figure}

\begin{figure}[H]
  \centering
  \includegraphics[width=0.62\textwidth]{v15_dotproduct.png}
  \caption{The two vectors of the formula shown together: the
  gradient $\nabla f(x(c),y(c))$ (crimson) and the velocity
  $(x'(c),y'(c))$ (blue), with the unit vector $\hat v$ along the
  velocity direction shown dashed.  The thick purple segment is
  \emph{not} a third vector: it is the scalar projection
  $\nabla f\cdot\hat v$, a length measured along the $\hat v$-axis
  (negative here, since the segment points opposite to $\hat v$),
  obtained by dropping a perpendicular (dotted grey) from the tip
  of $\nabla f$ onto the line through $\hat v$.  The quantity we
  actually want, $h'(c)=\nabla f\cdot(x',y')$, is a different
  number still --- the dot product with the \emph{un-normalised}
  velocity vector, $|(x',y')|$ times the scalar projection shown ---
  and is reported separately, in the text box, precisely to avoid
  drawing a number as though it were a vector.  In general,
  $\nabla f\cdot v = |\nabla f|\,|v|\cos\alpha$, where $\alpha$ is
  the angle between the two vectors; when they are perpendicular,
  $h'(c)=0$, matching the fact that gradients are orthogonal to
  level curves.}
  \label{fig:vova-env-15}
\end{figure}

\paragraph{Specialising to the envelope theorem: $y(c)=c$.}
The entire series above was deliberately general, with $x(c)$ and
$y(c)$ both free functions of $c$.  The connection to
\S\ref{subsec:envelope_simple}--\S\ref{subsec:envelope_operator} is
the special case where the \emph{second} variable is the identity:
$y(c) \equiv c$, so $y'(c)\equiv 1$ for every $c$.  Relabel
$x(c) \to x^*(c)$ (the optimiser) and write $f(x,c)$ in place of
$f(x,y)$.  The boxed chain-rule formula above becomes
\[
  \frac{d}{dc}\,f\big(x^*(c),\,c\big)
  = \frac{\partial f}{\partial x}\big(x^*(c),c\big)\cdot
    \underbrace{\frac{dx^*}{dc}}_{=\,x'(c)}
  \;+\;
  \frac{\partial f}{\partial c}\big(x^*(c),c\big)\cdot
  \underbrace{1}_{=\,y'(c)},
\]
which is exactly the chain-rule expansion used to prove the
elementary envelope theorem.  The first term vanishes by the
first-order condition $\partial f/\partial x\,(x^*(c),c)=0$, leaving
$V'(c) = \partial f/\partial c\,(x^*(c),c)$, the boxed result of
\S\ref{subsec:envelope_simple}.  In the dot-product language of
Figure~\ref{fig:vova-env-15}: the velocity vector in this special
case is $(x'(c),\,1)$ --- always a vector with second component
exactly $1$ --- and the vanishing of the first term is the
geometric statement that the \emph{first component} of $\nabla f$
(the $\partial f/\partial x$ direction) contributes nothing to the
dot product at the optimiser, because that component of $\nabla f$
is itself zero there.  The full generality of the picture series
was needed to see clearly that this vanishing is special to the
optimiser, and not a property of the dot-product formula itself,
which in general has both terms contributing.

%--------------------------------------------------------------------
\subsection{The Probabilist's continuation: the dot product as a function of
$c$, drawn on common abstract axes}
\label{subsec:envelope_uv_abstract}
%--------------------------------------------------------------------

The picture series above (\S\ref{subsec:envelope_chainrule_pictures})
reached the dot product $\nabla f\cdot(x',y')$ by drawing the two
vectors in their own native planes and resorting to a projection
construction to depict the resulting number.  The Probabilist's continuation,
from the same morning, reaches the identical place by a more direct
route: put both vectors in one common, abstract coordinate system
from the start, and treat the dot product not as a single
projection diagram but as an ordinary function of $c$.

\paragraph{Step 1: common abstract axes $(a,b)$.}
Rather than keeping the velocity vector in an $(x',y')$-plane and
the gradient in a separately-labelled plane, relabel both as living
in one shared abstract plane with axes $(a,b)$ --- emphasising that
\emph{both} vectors are just ordered pairs of real numbers, and the
labels $x',y'$ versus $\partial_xf,\partial_yf$ are bookkeeping, not
a difference in the kind of object.

\begin{figure}[H]
  \centering
  \includegraphics[width=0.62\textwidth]{v16_ab_axes.png}
  \caption{Both parametric curves placed in one common $(a,b)$-plane.
  Blue: $u(c)=(a,b)=(x'(c),\,y'(c))$, the velocity vector, traced as
  $c$ varies.  Green: $v(c)=(a,b)=\big(\partial_xf(x(c),y(c)),\,
  \partial_yf(x(c),y(c))\big)$, the gradient vector, also traced as
  $c$ varies.  Putting both in the same $(a,b)$-plane (rather than
  in two visually separate planes, as in
  Figures~\ref{fig:vova-env-8} and \ref{fig:vova-env-15}) makes
  plain that $u$ and $v$ are directly comparable objects of the same
  type, both living in $\mathbb{R}^2$.}
  \label{fig:vova-env-16}
\end{figure}

\paragraph{Step 2: fix one $c=c_0$ and isolate the two vectors.}
Pin down a single parameter value $c_0$ and draw only the two
resulting vectors $u(c_0)$ and $v(c_0)$, dropping the rest of each
curve.  Crucially, the dot product $u(c_0)\cdot v(c_0)$ is then
shown not as a geometric arrow but as a labelled point on a separate
number line below --- directly resolving the earlier confusion
(\S\ref{subsec:envelope_chainrule_pictures}, on the corrected
Figure~\ref{fig:vova-env-15}) about whether the dot product should
be drawn as a vector at all.  It should not: it is a single real
number, and a number line is the correct picture for it.

\begin{figure}[H]
  \centering
  \includegraphics[width=0.5\textwidth]{v17_uv_single_c.png}
  \caption{For one fixed $c=c_0$: the vector $u$ (blue) and the
  vector $v$ (green) in the $(a,b)$-plane.  Below, a separate number
  line carries the single point $u\cdot v$ --- the dot product,
  drawn exactly as what it is: a scalar, located on
  $\mathbb{R}$, not an arrow in $\mathbb{R}^2$.}
  \label{fig:vova-env-17}
\end{figure}

\paragraph{Step 3: let $c$ vary --- the dot product becomes a
function $\mathbb{R}\to\mathbb{R}$.}
The dot product itself, $(\cdot\,,\cdot):\mathbb{R}^2\times
\mathbb{R}^2\to\mathbb{R}$, is a fixed bilinear functional --- it
does not depend on $c$.  What depends on $c$ are its two arguments,
$u(c)$ and $v(c)$.  Composing,
\[
  c \;\longmapsto\; u(c)\cdot v(c) \;=\; h'(c)
\]
is an ordinary function from $\mathbb{R}$ to $\mathbb{R}$, exactly
like any other univariate function --- and so it can be drawn as an
ordinary 2-d graph over the $c$-axis, with no vectors or planes
needed anymore.

\begin{figure}[H]
  \centering
  \includegraphics[width=0.5\textwidth]{v18_hprime_of_c.png}
  \caption{The function $c\mapsto h'(c) = u(c)\cdot v(c)$, plotted
  directly as a curve over the $c$-axis.  All of the vector geometry
  of Figures~\ref{fig:vova-env-16}--\ref{fig:vova-env-17} has been
  absorbed into evaluating the dot product at each $c$; the result
  is just an ordinary real-valued function of one real variable,
  no different in kind from $x(c)$ or $y(c)$ themselves.}
  \label{fig:vova-env-18}
\end{figure}

\paragraph{Why this route is cleaner.}
The projection construction of
\S\ref{subsec:envelope_chainrule_pictures}
(Figure~\ref{fig:vova-env-15}) is geometrically informative --- it
shows \emph{why} the dot product has the value it does, via
$|\nabla f|\,|v|\cos\alpha$ --- but it requires care to avoid
mis-drawing the resulting scalar as a vector, exactly the error
flagged earlier in this section.  The Probabilist's route sidesteps that risk
entirely: by drawing $u\cdot v$ on a number line at a single $c$
(Figure~\ref{fig:vova-env-17}), and then immediately treating $h'(c)$
as an ordinary function of $c$ (Figure~\ref{fig:vova-env-18}), there
is no step at which a scalar quantity is ever rendered with an
arrowhead. The two approaches are complementary: the projection
picture explains the \emph{value} of the dot product geometrically;
this picture keeps the \emph{type} of the dot product (a number,
varying with $c$) unambiguous throughout.


%--------------------------------------------------------------------
\subsection{Three problems and the triangle connecting them}
\label{subsec:envelope_lm_summary}
%--------------------------------------------------------------------

We state three problems separately, in their own standard notation,
and then connect them in pairs.

\paragraph{Problem 1: constrained optimisation, parametrised by a
right-hand side $c$.}
Maximise $f(x)$ subject to $g(x)=c$, where $c$ is a parameter of the
constraint's right-hand side. The Lagrangian is
\[
  \mathcal{L}(x,\lambda;c) = f(x) - \lambda\big(g(x)-c\big).
\]
For each $c$, a stationary point $\big(x^*(c),\lambda^*(c)\big)$
solves $\nabla_{x,\lambda}\mathcal{L}=0$. Define the \emph{value
function}
\[
  V(c) := \mathcal{L}\big(x^*(c),\lambda^*(c);\,c\big)
        = f\big(x^*(c)\big),
\]
a function of $c$ alone. It can be differentiated in $c$; we will
need $V'(c)$ below.

\begin{remark}
This is the single-variable Lagrangian, matching Part~5b-a
throughout. It is deliberately \emph{not} the two-variable $f(x,y)$
of the chain-rule picture series
(\S\ref{subsec:envelope_chainrule_pictures}), where a genuine
second variable $y$ appears alongside a path $x(c),y(c)$. Side~A
below needs only one variable, so we state Problem~1 with one
variable from the start.
\end{remark}

\paragraph{Problem 2: the envelope theorem.}
In the simplest setting --- $\max_x f(x,c)$, unconstrained, no $g$
--- write $x^*(c)=\arg\max_xf(x,c)$ and $V(c)=f(x^*(c),c)$. The
envelope theorem states
\[
  V'(c) = \frac{\partial f}{\partial c}\Big(x^*(c),c\Big),
\]
proved by the chain rule: the term $\partial f/\partial
x\cdot dx^*/dc$ vanishes because $\partial f/\partial x=0$ at
$x^*(c)$. The same cancellation, applied jointly in
$(x,y,\lambda)$, gives Problem~1's value function
\[
  V'(c) = \frac{\partial\mathcal{L}}{\partial c}
  \bigg|_{x^*(c),y^*(c),\lambda^*(c)} = \lambda^*(c),
\]
since $\partial\mathcal{L}/\partial c=\lambda$. This is the
\emph{shadow price} reading: $\lambda^*$ is the rate of change of
the optimised value as the constraint level moves.

\paragraph{Problem 3: the score test, informally.}
Take a log-likelihood $\ell(\theta)$ and a simple null
$H_0:\theta=\theta_0$. The score is $s(\theta)=\ell'(\theta)$. Under
$H_0$, the identity $\E_{\theta_0}[s(\theta_0,X)]=0$ holds, and by
the CLT the sample average score concentrates around zero and
shrinks at rate $1/\sqrt{n}$ as the sample size grows. The test
rejects $H_0$ when the observed score $s(\theta_0)$ is too far from
zero relative to this shrinking spread --- formalised as
$\xi_{LM}=s(\theta_0)^2/\mathcal{I}(\theta_0)\xrightarrow{d}\chi^2_1$.

\paragraph{Side A: Problem 1 $\leftrightarrow$ Problem 2.}
Apply Problem~1 with $x\to\theta$, $f(x)\to\ell(\theta)$, and the
constraint $g(\theta)=\theta=c$ (so $g$ is the identity). The
genuine difficulty: the feasible set $\{\theta=c\}$ is one point, so
$\theta^*(c)=c$ trivially, with no first-order condition in $\theta$
alone --- the naive match $x\to\theta$ fails on its own, since there
is nothing to differentiate. The fix is to apply Problem~2's theorem
not to $\theta$ alone but to the pair $(\theta,\lambda)$ jointly,
via $\mathcal{L}(\theta,\lambda;c)=\ell(\theta)-\lambda(\theta-c)$:
both first-order conditions now hold genuinely
($\ell'(\theta)-\lambda=0$ recovers $\lambda$; $-(\theta-c)=0$
recovers $\theta=c$), even though $(\theta^*,\lambda^*)$ is a
saddle, not a maximum --- legitimate because the chain-rule proof of
Problem~2 only used the first-order condition, never concavity.
Carrying this out: $V(c)=\ell(\theta^*(c))-\lambda^*(c)\cdot 0=\ell(c)$
directly, so $V'(c)=\ell'(c)$; and via the envelope theorem
$V'(c)=\lambda^*(c)$. Equating:
\[
  \boxed{\ \lambda^*(c) = \ell'(c).\ }
\]

\paragraph{The Probabilist's strengthening: the Lagrangian equals $V$
identically, not just up to a constant.}
The Probabilist's notes of June~21, 2026 revisit Problem~2 once more, this time
following Wikipedia's statement of the constrained envelope theorem.
The upgrade relative to the further generalisation of
\S\ref{subsec:envelope_operator} (which let the constraint be
$g(\theta,c)=0$ but kept $f(\theta)$ free of $c$) is to let $f$
itself depend on $c$ as well, matching Problem~2's
$f(x,c)$:
\[
  \max_x f(x,c) \quad\text{subject to}\quad g(x,c)=0,
  \qquad
  \mathcal{L}(x,\lambda;c) = f(x,c) - \lambda\,g(x,c).
\]
For each fixed $c$, the Probabilist takes $\big(x^*(c),\lambda^*(c)\big)$ to be
the stationary point of $\mathcal{L}(x,\lambda)$ --- a function of
the \emph{two} variables $x,\lambda$ at that fixed $c$, exactly the
saddle-point picture recalled at the start of this notebook entry
--- and sets $V(c):=f(x^*(c),c)=\max_{g=0}f(x,c)$, as before. The
envelope theorem as already stated gives the derivative identity
\[
  V'(c) = \frac{\partial \mathcal{L}}{\partial c}
  \Big(x^*(c),\lambda^*(c),c\Big). \tag{1}
\]
The Probabilist's observation is that, because the constraint is an
\emph{equality} ($g=0$, not $g\le 0$), a strictly stronger statement
holds: not merely that the two sides of~(1) share a derivative, but
that the underlying functions of $c$ coincide \emph{exactly}, with
no leftover constant of integration:
\[
  \mathcal{L}\big(x^*(c),\lambda^*(c),c\big) = V(c)
  \qquad\text{for every } c. \tag{2}
\]
The proof is immediate from the definition of $\mathcal{L}$ together
with feasibility of the optimiser:
\[
  \mathcal{L}\big(x^*(c),\lambda^*(c),c\big)
  = f\big(x^*(c),c\big) - \lambda^*(c)\,g\big(x^*(c),c\big)
  = V(c) - \lambda^*(c)\cdot 0 = V(c),
\]
the last step using $g(x^*(c),c)=0$: the optimiser is feasible at
\emph{every} $c$, so the constraint term vanishes identically, not
just at one value of $c$.

Differentiating~(2) directly recovers~(1), so (2)~$\Rightarrow$~(1);
the converse need not hold, since~(1) alone only says the two sides
agree up to an additive constant in $c$, while~(2) pins that
constant at exactly zero. As the Probabilist puts it: the two functions are
not merely shifted copies of one another --- they are equal, the
shift being zero.

This is not new content so much as an unarticulated special case of
what Side~A already computed: the step
``$V(c)=\ell(\theta^*(c))-\lambda^*(c)\cdot 0=\ell(c)$'' above
\emph{is} identity~(2), for the particular instance where
$f(\theta)$ does not depend on $c$ and $g(\theta,c)=\theta-c$. The Probabilist's
note packages this as a general fact about equality-constrained
Lagrangians, rather than a coincidence of that one computation.

\emph{Caveat (the Probabilist).} Identity~(2) rests on $g(x^*(c),c)=0$ holding
exactly --- genuine equality. If the constraint is instead an
inequality, $g(x,c)\le 0$, only the derivative identity~(1) is
asserted in general; (2) may fail, even though complementary
slackness still gives $\lambda^*(c)\,g(x^*(c),c)=0$ pointwise at
each $c$. The Probabilist flags this contrast without resolving it further
here.

\paragraph{Side B: Problem 1 $\leftrightarrow$ Problem 3.}
Renaming $c\to\theta_0$, Side~A's boxed identity reads
$\lambda^*=\ell'(\theta_0)=s(\theta_0)$: the Lagrange multiplier at
the constrained optimum \emph{is} the score at the null. This is
exactly the quantity Problem~3 squares and divides by Fisher
information to build $\xi_{LM}$. So Problem~1's $\lambda^*$ is not
merely analogous to Problem~3's score --- it is the same number,
under two names, which is why the test is called both the
``score test'' and the ``LM test.''

\paragraph{Side C: Problem 2 $\leftrightarrow$ Problem 3.}
The shadow-price reading of Side~A ($\lambda^*=V'(\theta_0)$, the
sensitivity of the optimised log-likelihood to where we place the
null) gives the score a second interpretation, beyond ``slope of
$\ell$ at $\theta_0$'': it is \emph{how much the constrained optimum
would improve if the null were nudged}. A score near zero means the
optimised value is insensitive to nudging $\theta_0$ --- the null is
already where the unconstrained optimum would be, consistent with
$H_0$. A large score means nudging $\theta_0$ would improve the fit
substantially --- the data pull away from the null, exactly what
Problem~3's test is built to detect. The CLT shrinkage in Problem~3
($s(\theta_0)\to 0$ at rate $1/\sqrt n$ under $H_0$) is then read,
via Side~A, as: the shadow price of the constraint vanishes
asymptotically when the constraint is true.

\paragraph{What does not change.}
Problem~3's test statistic needs none of Problems~1--2 to be
computed --- as emphasised in Part~5b-a, \S2.7. The triangle above
is explanatory: it shows the same number, $\lambda^*=s(\theta_0)$,
arising from three different questions, and gives the score test's
central quantity a economic name (shadow price) that the bare
formula does not supply on its own.

%--------------------------------------------------------------------
\subsection{A historical note: which came first, score or
multiplier?}
\label{subsec:envelope_lm_history}
%--------------------------------------------------------------------

The triangle above raises a natural question: is the Lagrangian
apparatus actually needed to define or compute the score test, or
is it decoration? The answer, confirmed by the history, is that it
is decoration --- but the chronology is the reverse of what one
might guess.

\paragraph{Rao (1948) came first, with no Lagrangian at all.}
C.~R. Rao's 1948 paper introduces the score test directly: the
statistic is built from the score $s(\theta_0)$ and the Fisher
information $\mathcal{I}(\theta_0)$, with $\xi=s(\theta_0)^2/
\mathcal{I}(\theta_0)\xrightarrow{d}\chi^2_1$ under $H_0$. No
constrained optimisation, no Lagrangian, no multiplier appears in
this original formulation --- it is Problem~3 above, stated and
proved on its own terms.

\paragraph{Silvey (1959) came eleven years later, adding the
Lagrangian framing.}
Aitchison and Silvey (1958) and Silvey (1959) approached the same
problem from the constrained-MLE side: maximise the log-likelihood
subject to the null's restrictions, form the Lagrangian, and observe
that the resulting multiplier $\lambda^*$ is (asymptotically)
distributed the same way as Rao's score statistic. Silvey's 1959
paper is the source of the name ``Lagrange Multiplier test.'' So
the order is: \emph{score test first, Lagrangian re-derivation
second} --- not the reverse.

\paragraph{Correcting one detail, confirming the main point.}
This means the natural guess --- that the test was first defined
via $\lambda^*$ and only later ``unburdened'' into the score form
--- has the chronology backwards: Rao's original score form never
carried Lagrangian baggage to begin with; Silvey's Lagrangian
version was an \emph{additional}, later route to the same number,
not a simplification of an earlier, more complicated one. But the
substantive conclusion is exactly right: computing and justifying
the test requires only $s(\theta_0)$ and $\mathcal{I}(\theta_0)$,
both of which are obtained by differentiation alone, with no
optimisation or stationarity condition involved (Problem~3, on its
own, needs none of Problem~1's machinery). The Lagrangian route
(Side~A above) shows that $\lambda^*$ --- obtained by solving a
constrained-optimisation stationarity system --- happens to equal
$s(\theta_0)$, but as established earlier in this document, that
stationarity system is solving a problem whose feasible set is a
single point: there is no genuine optimisation content in it, and
the ``answer'' it returns is no different from what direct
differentiation already supplies.

\paragraph{Why both names survived, rather than one replacing the
other.}
There is no documented single event where the field renamed the
test from ``LM'' to ``score''; both names remain in active use
today, split roughly along disciplinary lines --- the score-test
name is more common in statistics (following Rao), the LM-test name
more common in econometrics (following Silvey, and popularised
further by Breusch and Pagan's 1980 use of it for the
heteroscedasticity test that motivated this entire document). The
coexistence reflects two independent discoveries of the same
asymptotic statistic, each retaining the name attached to its own
derivation route, rather than a deliberate later simplification of
one name into the other.

\end{document}
